\section{Code evolution across milestone versions}

% Requires in the main preamble:
% \usepackage{listings}
% \usepackage{xcolor}
%
% If not already defined in the report preamble, the following style block can
% remain here safely.
\definecolor{codebg}{RGB}{248,248,248}
\definecolor{codeframe}{RGB}{220,220,220}
\definecolor{codecomment}{RGB}{76,114,176}
\definecolor{codekeyword}{RGB}{180,60,40}
\definecolor{codestring}{RGB}{34,139,34}

\lstdefinestyle{appendixcode}{
  language=C++,
  basicstyle=\ttfamily\scriptsize,
  keywordstyle=\color{codekeyword}\bfseries,
  commentstyle=\color{codecomment},
  stringstyle=\color{codestring},
  numbers=left,
  numberstyle=\tiny,
  stepnumber=1,
  numbersep=8pt,
  showstringspaces=false,
  tabsize=4,
  breaklines=true,
  breakatwhitespace=false,
  columns=fullflexible,
  keepspaces=true,
  frame=single,
  rulecolor=\color{codeframe},
  backgroundcolor=\color{codebg},
  xleftmargin=1.2em,
  framexleftmargin=1.0em,
  aboveskip=0.8em,
  belowskip=0.8em,
  captionpos=b
}

This appendix presents the Cholesky implementation for each milestone version using the archived source snapshots stored in the repository. For v0--v3 the listings are taken from \texttt{history/versions/<tag>/src/cholesky.cpp}, while the v4 listing is taken from the current \texttt{src/cholesky.cpp}. For the OpenMP versions, short appendix-only explanatory comments are inserted around OpenMP-specific lines to make the parallel intent clear without rewriting project history.

\subsection{v0: Baseline serial kernel}

Tag: \texttt{cholesky\_v0\_baseline}.

This version is a direct serial implementation of the coursework baseline algorithm. It serves as the initial correctness and timing reference before any restructuring or parallelisation.

\begin{lstlisting}[style=appendixcode]
#include "mphil_dis_cholesky.h"

#include <chrono>
#include <cmath>
#include <cstddef>

double cholesky(double* c, int n) {
    if (c == nullptr || n < 0 || n > 100000) {
        return -1.0;
    }

    const std::size_t stride = static_cast<std::size_t>(n);
    const auto start = std::chrono::steady_clock::now();

    for (int p = 0; p < n; ++p) {
        const std::size_t p_index = static_cast<std::size_t>(p);
        const std::size_t diag_offset = p_index * stride + p_index;
        const double diag_entry = c[diag_offset];

        if (diag_entry <= 0.0) {
            return -1.0;
        }

        const double diag = std::sqrt(diag_entry);
        c[diag_offset] = diag;

        for (int j = p + 1; j < n; ++j) {
            const std::size_t j_offset = static_cast<std::size_t>(j);
            c[p_index * stride + j_offset] /= diag;
        }

        for (int i = p + 1; i < n; ++i) {
            const std::size_t i_offset = static_cast<std::size_t>(i);
            c[i_offset * stride + p_index] /= diag;
        }

        for (int j = p + 1; j < n; ++j) {
            const std::size_t j_offset = static_cast<std::size_t>(j);
            const double row_value = c[p_index * stride + j_offset];

            for (int i = p + 1; i < n; ++i) {
                const std::size_t i_offset = static_cast<std::size_t>(i);
                c[i_offset * stride + j_offset] -= c[i_offset * stride + p_index] * row_value;
            }
        }
    }

    const auto end = std::chrono::steady_clock::now();
    const std::chrono::duration<double> elapsed = end - start;
    return elapsed.count();
}
\end{lstlisting}

\subsection{v1: Serial kernel clean-up}

Tag: \texttt{cholesky\_v1\_serial\_opt}.

This version keeps the same mathematical algorithm as v0, but reduces loop overhead, reuses row pointers, and replaces repeated division with multiplication by a precomputed reciprocal.

\begin{lstlisting}[style=appendixcode]
#include "mphil_dis_cholesky.h"

#include <chrono>
#include <cmath>
#include <cstddef>

double cholesky(double* c, int n) {
    if (c == nullptr || n < 0 || n > 100000) {
        return -1.0;
    }

    const std::size_t stride = static_cast<std::size_t>(n);
    const auto start = std::chrono::steady_clock::now();

    for (int p = 0; p < n; ++p) {
        const std::size_t p_index = static_cast<std::size_t>(p);
        const std::size_t diag_offset = p_index * stride + p_index;
        double* const row_p = c + p_index * stride;
        const double diag_entry = c[diag_offset];

        if (diag_entry <= 0.0) {
            return -1.0;
        }

        const double diag = std::sqrt(diag_entry);
        const double inv_diag = 1.0 / diag;
        c[diag_offset] = diag;

        for (int j = p + 1; j < n; ++j) {
            const std::size_t j_offset = static_cast<std::size_t>(j);
            row_p[j_offset] *= inv_diag;
        }

        for (int i = p + 1; i < n; ++i) {
            const std::size_t i_offset = static_cast<std::size_t>(i);
            c[i_offset * stride + p_index] *= inv_diag;
        }

        for (int i = p + 1; i < n; ++i) {
            const std::size_t i_offset = static_cast<std::size_t>(i);
            double* const row_i = c + i_offset * stride;
            const double col_value = row_i[p_index];

            for (int j = p + 1; j < n; ++j) {
                const std::size_t j_offset = static_cast<std::size_t>(j);
                row_i[j_offset] -= col_value * row_p[j_offset];
            }
        }
    }

    const auto end = std::chrono::steady_clock::now();
    const std::chrono::duration<double> elapsed = end - start;
    return elapsed.count();
}
\end{lstlisting}

\subsection{v2: Blocked serial formulation}

Tag: \texttt{cholesky\_v2\_serial\_blocked}.

This version introduces panel factorisation, block-column solve, and tiled trailing-matrix updates. The purpose is to improve locality while remaining fully serial.

\begin{lstlisting}[style=appendixcode]
#include "mphil_dis_cholesky.h"

#include <chrono>
#include <cmath>
#include <cstddef>
#include <algorithm>
#include <cstdlib>

namespace {

int resolve_block_size() {
    constexpr int default_block_size = 32;
    const char* value = std::getenv("CHOLESKY_BLOCK_SIZE");
    if (value == nullptr) {
        return default_block_size;
    }

    const int parsed = std::atoi(value);
    return parsed > 0 ? parsed : default_block_size;
}

}  // namespace

double cholesky(double* c, int n) {
    if (c == nullptr || n < 0 || n > 100000) {
        return -1.0;
    }

    const std::size_t stride = static_cast<std::size_t>(n);
    const auto start = std::chrono::steady_clock::now();
    const int panel_block_size = resolve_block_size();
    const int trailing_block_size = panel_block_size;

    auto at = [c, stride](int row, int col) -> double& {
        return c[static_cast<std::size_t>(row) * stride + static_cast<std::size_t>(col)];
    };

    for (int block_start = 0; block_start < n; block_start += panel_block_size) {
        const int block_end = std::min(block_start + panel_block_size, n);

        // Factor the diagonal panel one column at a time while keeping it mirrored.
        for (int p = block_start; p < block_end; ++p) {
            const double diag_entry = at(p, p);
            if (diag_entry <= 0.0) {
                return -1.0;
            }

            const double diag = std::sqrt(diag_entry);
            const double inv_diag = 1.0 / diag;
            at(p, p) = diag;

            for (int i = p + 1; i < block_end; ++i) {
                const double value = at(i, p) * inv_diag;
                at(i, p) = value;
                at(p, i) = value;
            }

            for (int j = p + 1; j < block_end; ++j) {
                const double panel_value = at(j, p);
                for (int i = j; i < block_end; ++i) {
                    const double updated = at(i, j) - at(i, p) * panel_value;
                    at(i, j) = updated;
                    if (i != j) {
                        at(j, i) = updated;
                    }
                }
            }
        }

        // Solve the block column below the panel using the newly factorized block.
        for (int i = block_end; i < n; ++i) {
            for (int j = block_start; j < block_end; ++j) {
                double value = at(i, j);
                for (int p = block_start; p < j; ++p) {
                    value -= at(i, p) * at(j, p);
                }
                value /= at(j, j);
                at(i, j) = value;
                at(j, i) = value;
            }
        }

        // Update the trailing matrix with blocked rank-k style tiles.
        for (int i_block = block_end; i_block < n; i_block += trailing_block_size) {
            const int i_end = std::min(i_block + trailing_block_size, n);

            for (int j_block = block_end; j_block <= i_block; j_block += trailing_block_size) {
                const int j_end = std::min(j_block + trailing_block_size, n);

                for (int i = i_block; i < i_end; ++i) {
                    const int j_limit = std::min(j_end, i + 1);
                    for (int j = j_block; j < j_limit; ++j) {
                        double value = at(i, j);
                        for (int p = block_start; p < block_end; ++p) {
                            value -= at(i, p) * at(j, p);
                        }
                        at(i, j) = value;
                        if (i != j) {
                            at(j, i) = value;
                        }
                    }
                }
            }
        }
    }

    const auto end = std::chrono::steady_clock::now();
    const std::chrono::duration<double> elapsed = end - start;
    return elapsed.count();
}
\end{lstlisting}

\subsection{v3: First OpenMP version}

Tag: \texttt{cholesky\_v3\_openmp}.

This version keeps the blocked structure of v2 and introduces OpenMP in the trailing update. That region is the natural first parallel hotspot because tiles are independent for a fixed panel.

\begin{lstlisting}[style=appendixcode]
#include "mphil_dis_cholesky.h"

#include <chrono>
#include <cmath>
#include <cstddef>
#include <algorithm>
#include <cstdlib>
#include <omp.h>  // OpenMP is used for the first parallel version in the trailing-update hotspot.

namespace {

int resolve_block_size() {
    constexpr int default_block_size = 32;
    const char* value = std::getenv("CHOLESKY_BLOCK_SIZE");
    if (value == nullptr) {
        return default_block_size;
    }

    const int parsed = std::atoi(value);
    return parsed > 0 ? parsed : default_block_size;
}

}  // namespace

double cholesky(double* c, int n) {
    if (c == nullptr || n < 0 || n > 100000) {
        return -1.0;
    }

    const std::size_t stride = static_cast<std::size_t>(n);
    const auto start = std::chrono::steady_clock::now();
    const int panel_block_size = resolve_block_size();
    const int trailing_block_size = panel_block_size;

    auto at = [c, stride](int row, int col) -> double& {
        return c[static_cast<std::size_t>(row) * stride + static_cast<std::size_t>(col)];
    };

    for (int block_start = 0; block_start < n; block_start += panel_block_size) {
        const int block_end = std::min(block_start + panel_block_size, n);

        // Factor the diagonal panel one column at a time while keeping it mirrored.
        for (int p = block_start; p < block_end; ++p) {
            const double diag_entry = at(p, p);
            if (diag_entry <= 0.0) {
                return -1.0;
            }

            const double diag = std::sqrt(diag_entry);
            const double inv_diag = 1.0 / diag;
            at(p, p) = diag;

            for (int i = p + 1; i < block_end; ++i) {
                const double value = at(i, p) * inv_diag;
                at(i, p) = value;
                at(p, i) = value;
            }

            for (int j = p + 1; j < block_end; ++j) {
                const double panel_value = at(j, p);
                for (int i = j; i < block_end; ++i) {
                    const double updated = at(i, j) - at(i, p) * panel_value;
                    at(i, j) = updated;
                    if (i != j) {
                        at(j, i) = updated;
                    }
                }
            }
        }

        // Solve the block column below the panel using the newly factorized block.
        for (int i = block_end; i < n; ++i) {
            for (int j = block_start; j < block_end; ++j) {
                double value = at(i, j);
                for (int p = block_start; p < j; ++p) {
                    value -= at(i, p) * at(j, p);
                }
                value /= at(j, j);
                at(i, j) = value;
                at(j, i) = value;
            }
        }

        // Update the trailing matrix with blocked rank-k style tiles.
        // OpenMP is introduced here first because these trailing tiles are independent for a fixed panel and dominate the runtime.
// Appendix note: This OpenMP parallel-for distributes independent work across threads in the main parallel region.
#pragma omp parallel for schedule(static) if (n - block_end > trailing_block_size)
        for (int i_block = block_end; i_block < n; i_block += trailing_block_size) {
            const int i_end = std::min(i_block + trailing_block_size, n);

            for (int j_block = block_end; j_block <= i_block; j_block += trailing_block_size) {
                const int j_end = std::min(j_block + trailing_block_size, n);

                for (int i = i_block; i < i_end; ++i) {
                    const int j_limit = std::min(j_end, i + 1);
                    for (int j = j_block; j < j_limit; ++j) {
                        double value = at(i, j);
                        for (int p = block_start; p < block_end; ++p) {
                            value -= at(i, p) * at(j, p);
                        }
                        at(i, j) = value;
                        if (i != j) {
                            at(j, i) = value;
                        }
                    }
                }
            }
        }
    }

    const auto end = std::chrono::steady_clock::now();
    const std::chrono::duration<double> elapsed = end - start;
    return elapsed.count();
}
\end{lstlisting}

\subsection{v4: Final tuned OpenMP kernel}

Tag: \texttt{cholesky\_v4\_parallel\_opt}.

This is the final accepted implementation. It extends the OpenMP version with runtime-configurable scheduling, parallel block-column solve, and SIMD-assisted inner kernels.

\begin{lstlisting}[style=appendixcode]
#include "mphil_dis_cholesky.h"

#include <algorithm>
#include <chrono>
#include <cmath>
#include <cstddef>
#include <cstdlib>
#include <omp.h>  // OpenMP drives the tuned parallel trailing-update kernel.
#include <string>

namespace {

// Reads a positive integer override from the environment and falls back to a
// known-safe default when the variable is unset or invalid.
int resolve_positive_env_or_default(const char* name, int default_value) {
    const char* value = std::getenv(name);
    if (value == nullptr) {
        return default_value;
    }

    const int parsed = std::atoi(value);
    return parsed > 0 ? parsed : default_value;
}

// Returns the blocked-kernel tile size used by the current implementation.
int resolve_block_size() {
    return resolve_positive_env_or_default("CHOLESKY_BLOCK_SIZE", 32);
}

// Stores the OpenMP runtime schedule policy selected for the current run.
// Appendix note: This helper stores the runtime OpenMP schedule policy and chunk size used by the tuned final kernel.
struct OpenMPScheduleConfig {
// Appendix note: The schedule kind is kept explicit so the benchmarking workflow can compare static, dynamic, and guided policies.
    omp_sched_t kind = omp_sched_guided;
    int chunk = 1;
};

// Resolves the OpenMP runtime schedule configuration used by the v4 kernel.
OpenMPScheduleConfig resolve_openmp_schedule() {
    OpenMPScheduleConfig config;

    if (const char* schedule_value = std::getenv("CHOLESKY_OMP_SCHEDULE")) {
        const std::string schedule(schedule_value);
        if (schedule == "dynamic") {
            config.kind = omp_sched_dynamic;
        } else if (schedule == "guided") {
            config.kind = omp_sched_guided;
        } else if (schedule == "auto") {
            config.kind = omp_sched_auto;
        } else {
            config.kind = omp_sched_static;
        }
    }

    if (const char* chunk_value = std::getenv("CHOLESKY_OMP_CHUNK")) {
        const int parsed_chunk = std::atoi(chunk_value);
        if (parsed_chunk > 0) {
            config.chunk = parsed_chunk;
        }
    }

    return config;
}

}  // namespace

// Computes an in-place Cholesky factorization of a symmetric positive-definite
// row-major matrix and returns the elapsed wall-clock time in seconds.
double cholesky(double* c, int n) {
    if (c == nullptr || n < 0 || n > 100000) {
        return -1.0;
    }

    const std::size_t stride = static_cast<std::size_t>(n);
    const int block_size = resolve_block_size();
    const OpenMPScheduleConfig openmp_schedule = resolve_openmp_schedule();

    // Apply the runtime schedule selected for the current run.
// Appendix note: The OpenMP runtime is configured here so benchmark runs can tune schedule policy without recompiling.
    omp_set_schedule(openmp_schedule.kind, openmp_schedule.chunk);

    auto row_ptr = [c, stride](int row) -> double* {
        return c + static_cast<std::size_t>(row) * stride;
    };

    const auto start = std::chrono::steady_clock::now();

    for (int block_start = 0; block_start < n; block_start += block_size) {
        const int block_end = std::min(block_start + block_size, n);

        // Factor the diagonal panel and preserve the mirrored upper triangle.
        for (int p = block_start; p < block_end; ++p) {
            double* const row_p = row_ptr(p);
            const double diag_entry = row_p[p];
            if (diag_entry <= 0.0) {
                return -1.0;
            }

            const double diag = std::sqrt(diag_entry);
            const double inv_diag = 1.0 / diag;
            row_p[p] = diag;

            for (int i = p + 1; i < block_end; ++i) {
                double* const row_i = row_ptr(i);
                row_i[p] *= inv_diag;
                row_p[i] = row_i[p];
            }

            for (int j = p + 1; j < block_end; ++j) {
                double* const row_j = row_ptr(j);
                const double panel_value = row_j[p];

                for (int i = j; i < block_end; ++i) {
                    double* const row_i = row_ptr(i);
                    const double updated = row_i[j] - row_i[p] * panel_value;
                    row_i[j] = updated;
                    if (i != j) {
                        row_j[i] = updated;
                    }
                }
            }
        }

        // Rows below the panel are independent once the panel is factored.
        // Parallelize over rows to solve the block column.
// Appendix note: This OpenMP parallel-for distributes independent work across threads in the main parallel region.
#pragma omp parallel for schedule(runtime) if (n - block_end > block_size)
        for (int i = block_end; i < n; ++i) {
            double* const row_i = row_ptr(i);

            for (int j = block_start; j < block_end; ++j) {
                double* const row_j = row_ptr(j);
                double dot = 0.0;

                // Vectorize the panel dot product reduction.
// Appendix note: This OpenMP SIMD reduction asks the compiler to vectorise the inner dot product safely.
 #pragma omp simd reduction(+:dot)
                for (int p = block_start; p < j; ++p) {
                    dot += row_i[p] * row_j[p];
                }

                row_i[j] = (row_i[j] - dot) / row_j[j];
                row_j[i] = row_i[j];
            }
        }

        // Each outer tile owns disjoint lower-triangular rows in the trailing matrix.
        // Parallelize over those tiles in the main OpenMP hotspot.
// Appendix note: This OpenMP parallel-for distributes independent work across threads in the main parallel region.
#pragma omp parallel for schedule(runtime) if (n - block_end > block_size)
        for (int i_block = block_end; i_block < n; i_block += block_size) {
            const int i_end = std::min(i_block + block_size, n);

            for (int j_block = block_end; j_block <= i_block; j_block += block_size) {
                const int j_end = std::min(j_block + block_size, n);

                for (int i = i_block; i < i_end; ++i) {
                    double* const row_i = row_ptr(i);
                    const int j_limit = std::min(j_end, i + 1);

                    for (int j = j_block; j < j_limit; ++j) {
                        double* const row_j = row_ptr(j);
                        double dot = 0.0;

                        // Vectorize the rank-k accumulation inside each tile.
// Appendix note: This OpenMP SIMD reduction asks the compiler to vectorise the inner dot product safely.
 #pragma omp simd reduction(+:dot)
                        for (int p = block_start; p < block_end; ++p) {
                            dot += row_i[p] * row_j[p];
                        }

                        row_i[j] -= dot;
                        if (i != j) {
                            row_j[i] = row_i[j];
                        }
                    }
                }
            }
        }
    }

    const auto end = std::chrono::steady_clock::now();
    const std::chrono::duration<double> elapsed = end - start;
    return elapsed.count();
}
\end{lstlisting}

